\section{The strong-coupling gap}\label{sec:gap}

This section states the central result of the record, from the fifth-reference manuscript \texttt{FIFTH\_REFERENCE.md} (equations F1--F42; reproduced as Appendix~A of the fifth-reference reader), with its proof. It was first stated in the reader of four workbenches \cite{WorkbenchReader}, where it had three referee passes (Appendix~\ref{app:verification}).

\subsection{The statement}

\begin{theorem}[the fifth-reference gap bound; F26--F42]\label{thm:ym01}
Let $(m_i,t_i)$, $i=1,\dots,5$, be the ten rational numbers of the workbench's F26, $\ell_5(x)=\sum(m_i+4t_i)x^i$, $\delta_5(x)=\sum_{i,j\le5,\ i+j\ge6}3(m_it_j+m_jt_i)x^{i+j}$, $\mathcal D_5=(1-\ell_5)^2-\frac83\delta_5$, $\alpha_5$ the first positive root of $\mathcal D_5$, and
\[
d_5(x)=\tfrac32\bigl(1+\sqrt{\mathcal D_5(x)}\bigr)+\textstyle\sum_{i=1}^5\bigl(\tfrac32m_i-6t_i\bigr)x^i.
\]
If the $(m_i,t_i)$ bound the local norms of the vacuum coefficients as the workbench's computations assert, then for every $L\ge2$, every $a>0$ and every $0<\xi\le\alpha_5$,
\[
\Delta_L\ge\kappa\,d_5(\xi).
\]
In the coupling,
\[
\xi\le\alpha_5\iff g^2\ge\frac1{2\sqrt{\alpha_5}}=3.6835519839857273\ldots,\qquad\alpha_5=0.0184249535761176166818\ldots;
\]
$d_5$ decreases on $[0,\alpha_5]$ from $d_5(0)=3$ (at $\xi=0$, $\Delta_L=3\kappa$) to $d_5(\alpha_5)=1.5453\ldots$; so $\Delta_L\ge1.5453\,\kappa$ for all $g^2\ge3.6836$, and, as $d_5(1/64)=1.9068301567\ldots$, $\Delta_L>1.9068\,\kappa$ for all $g^2\ge4$, uniformly in the box and the spacing. Equivalently, $q_{H_L-E_0}(\psi f)\ge\kappa\,d_5(\xi)\Var_{\psi^2}(f)$ for $f$ in the physical form domain, $\psi$ the normalised ground state (the workbench's form, F37).
\end{theorem}
\status{\textsf{Conditional} on the computed coefficient bounds of orders 3, 4 and 5 for the threshold $3.6836$ and the constants $d_5$ (the audit's inventory label is YM-01); its qualitative content, a gap on the whole physical space, uniform in the box and the spacing, above an explicit coupling, holds with audited inputs for $g^2\ge3.9781$ (Corollary~\ref{cor:ym-orderN}(b)).
\begin{itemize}
\item \emph{Order 1} is verified by hand.
\item \emph{Order 2} was recomputed independently in the first pass from explicit $SU(2)$ coefficient tensors. The exact values $m(v_2)=134.0673\ldots$ and $t(v_2)=19.7953\ldots$ lie below the workbench's $5834/39$ and $137/6$.
\item \emph{Order 5} (662 representatives, 124{,}864 anchored multisets, 5{,}726 rational primal/dual certificates) was reproduced by re-running the workbench's producer and auditor, with byte-identical output. The workbench's execution record reports these runs; its publication intake matched their hashes but did not repeat them (\texttt{INTAKE\_REVIEW.md}). The logic of order 5 was read, not re-implemented.
\item \emph{Orders 3--4} are inherited from earlier continuations and pinned by hash.
\end{itemize} The analytic argument is audited. $\alpha_5$, the threshold and $d_5$ were recomputed by exact bisection and by root isolation, and the monotonicity of $d_5$ on a $401$-point grid.}

\subsection{The proof: two lemmas and the analytic core}

\begin{lemma}[gauge-invariance Casimir inequality]\label{lem:casimir}
Let the links form a finite graph without loops (parallel links allowed) in which every cycle has length at least $g$; the cubic box has $g=4$. On every Peter--Weyl block $j$ that contains a nonzero gauge-invariant vector, $\sum_fj_f\ge g\,j_e$ for every link $e$. Hence $c_j\ge\frac{3g}2j_e$ for every $e$, $\sum_ej_e\le\frac23c_j$, and $c_j\ge\frac{3g}4$ on the nonconstant physical space; for the box, $c_j\ge6j_e$ and $c_j\ge3$. The first inequality holds for all real weights $j_f\ge0$ that satisfy the vertex inequalities $j_f\le\sum_{f'\ni v,\,f'\ne f}j_{f'}$, and it is an equality for $j=\frac12$ on a cycle of length $g$ through $e$ and $0$ elsewhere.
\end{lemma}
\status{Generalised here: the workbench's F7 (audited) is the case $g=4$ on the cubic lattice, proved there from the absence of triangles; the sharp constant on any graph is its girth. Checked exhaustively over half-integer spins with vertex singlets on the cube graph ($1{,}012$ nonzero assignments with every $j_e\le1$; constant $4$), $K_{3,3}$ ($4$), the $5$-cycle and the Petersen graph ($5$), the $6$-cycle ($6$) and a triangle ($3$); and by linear programming over real weights with the vertex inequalities only, on $13$ graphs of girth $3$ to $8$ (among them the Petersen, Heawood, McGee and Tutte--Coxeter graphs, the dodecahedron and the box $\{-1,0,1\}^3$), on three parallel links ($2$), and on the bridge between two triangles. On every link that lies on a shortest cycle the optimum is the girth; on the bridge, which lies on no cycle, it is $4$, larger than the girth $3$. So the girth is the best constant uniform over the links; Remark~\ref{rem:perlink} gives the exact constant for each link.}
\begin{proof}
A block contains a nonzero invariant exactly when, at every vertex $v$, the tensor product of the representations of spins $j_f$ ($f\ni v$) contains the trivial representation; by the Clebsch--Gordan rule this gives the vertex inequalities, which are all that is used.
\emph{A matching at each vertex.} At $v$, the inequalities say that no $j_f$ exceeds half of $J_v=\sum_{f\ni v}j_f$. That is the condition for a transport plan between the links at $v$, with supplies and demands $j_f$ and the diagonal forbidden, by Gale's supply--demand theorem: a set of two or more columns can be served by every row, and a single column $f$ by the rows $f'\ne f$, which carry $J_v-j_f\ge j_f$. Symmetrising the plan gives a symmetric $M^v\ge0$ with zero diagonal and row sums $j_f$.
\emph{A circulation.} Give each orientation (dart) of a link $f$ the weight $j_f/2$, and let the dart arriving at $v$ along $f$ pass to the dart leaving $v$ along $f'\ne f$ with weight $M^v_{ff'}/2$. The row and column sums of $M^v$ make this a circulation on the darts, which is therefore a nonnegative combination $\sum_Cw_C[C]$ of cyclic dart sequences $C$, that is, closed walks in which consecutive links (cyclically) differ. Counting weights, $\sum_Cw_C|C|=\sum_fj_f$ and $\sum_Cw_Cn_e(C)=j_e$, where $n_e(C)$ is the number of traversals of $e$ by $C$.
\emph{Each walk.} Cut a walk $C$ with $n_e(C)\ge1$ at its traversals of $e=uw$. Each piece avoids $e$ and is nonempty, since two traversals of $e$ are never consecutive and $u\ne w$. A piece that ends at the endpoint of $e$ where it started is a closed walk in which consecutive links differ, so it contains a cycle and has length at least $g$; a piece joining $u$ and $w$ contains a path, which with $e$ forms a cycle, so it has length at least $g-1$. Hence $|C|\ge g\,n_e(C)$, and summing, $\sum_fj_f\ge g\,j_e$.
\emph{The Casimir.} Every nonzero spin is at least $\frac12$, so $j_f(j_f+1)\ge\frac32j_f$; hence $c_j\ge\frac32\sum_fj_f$, which is $\sum_fj_f\le\frac23c_j$, and $c_j\ge\frac{3g}2j_e$. On a nonconstant physical block some $j_e\ge\frac12$, so $c_j\ge\frac{3g}4$. For the equality case, $j=\frac12$ on a cycle of length $g$ satisfies every vertex inequality and has $\sum_fj_f=\frac g2=g\,j_e$.
\end{proof}

\begin{remark}[the exact constant for each link]\label{rem:perlink}
For a link $e=uw$ let $g_e$ be the length of a shortest cycle through $e$ ($g_e=\infty$ if there is none), and $\ell_x$ the length of a shortest closed walk at $x$ in the graph without $e$ in which consecutive links differ, except possibly the last and the first. Then $\sum_fj_f\ge\gamma_e\,j_e$ with $\gamma_e=\min\bigl(g_e,1+\frac12(\ell_u+\ell_w)\bigr)\ge g$, and $\gamma_e$ is attained by half-integer spins with vertex singlets.
\end{remark}
\status{Proved here; found by the referee pass of version~1. The referee's linear programs give $\gamma_e$ on every link of ten graphs, including both bridge examples.}
\begin{proof}
Cut each walk $C$ of the circulation at its traversals of $e$, as in the proof of Lemma~\ref{lem:casimir}. A piece joining $u$ and $w$ has length at least $g_e-1$. A piece that returns to its starting endpoint $x$ is a closed walk at $x$ in the graph without $e$, of length at least $\ell_x$. After a piece that returns to $w$ the next traversal of $e$ runs from $w$ to $u$, and conversely; so, as $C$ is closed, the returning pieces come in pairs, one at $u$ and one at $w$. Hence $|C|\ge n_e(C)\gamma_e$, and summing, $\sum_fj_f\ge\gamma_ej_e$.
For attainment, spin $\frac12$ on a shortest cycle through $e$ gives the ratio $g_e$. For the other value, take $C=eP_we^{-1}P_u$, with $P_w,P_u$ shortest walks of the kind defining $\ell_w,\ell_u$, and put $j_f=\frac12n_f(C)$. Consecutive links of $C$ differ, which gives the vertex inequalities. The spins at a vertex sum to its number of visits by $C$, an integer, so a singlet exists. The ratio is $|C|/n_e(C)=1+\frac12(\ell_u+\ell_w)$.
\end{proof}

\begin{lemma}[ground-state transform]\label{lem:gst}
Let $\Gamma(f,h)=\sum_{e,a}(X_{e,a}f)(X_{e,a}h)$, $P_H$ the Haar mean and $Q_H=I-P_H$. For real or complex $v$, $e^v$ is an eigenfunction of $H_L$ with eigenvalue $E$ iff $Kv=\xi S+Q_H\Gamma(v,v)$ and $E=\kappa(2M_L\xi-P_H\Gamma(v,v))$; then $(H_L-E)(e^vh)=\kappa e^v(Kh-2\Gamma(v,h))$ for every $h$.
\end{lemma}
\status{Proved here; the workbench's F33--F36, audited. The identities are algebraic, so they hold for complex $v$ and complex $\xi$; only the positivity of $e^v$, used for the ground state, needs $v$ real.}
\begin{proof}
From $X^2(e^vh)=e^v(h(Xv)^2+2(Xv)(Xh)+hX^2v+X^2h)$ one gets $K(e^vh)=e^v\bigl(h(Kv-\Gamma(v,v))-2\Gamma(v,h)+Kh\bigr)$. With $h=1$, $H_Le^v=\kappa e^v(Kv-\Gamma(v,v)+\xi(2M_L-S))$, which equals $Ee^v$ exactly when $Kv-\Gamma(v,v)-\xi S$ is the constant $E/\kappa-2M_L\xi$; taking Haar means ($P_HKv=0$, $P_HS=0$) identifies the constant as $-P_H\Gamma(v,v)$. Inserting this into the formula for $K(e^vh)$ gives the last identity.
\end{proof}

The ground state is unique and positive, so a solution $v$ of the vacuum equation gives the ground state $e^v/\|e^v\|$. With $v_1=S/3$ (each $W_p$ has Casimir $3$) and $B=K^{-1}Q_H\Gamma$ the equation is $v=\xi v_1+B(v,v)$, whose formal solution is $\sum\xi^nv_n$ with $v_n=\sum_{i+j=n}B(v_i,v_j)$; this exponential (coupled-cluster) ansatz is the one the workbench credits, for Hamiltonian lattice field theory, to Schütte, Zheng and Hamer \cite{SchutteZhengHamer}.

\emph{The analytic core} (the workbench's F4--F32; audited). In the Fourier-algebra norm $\|f\|_X=\sum_j\|A_j(f)\|_1$, which is submultiplicative (Eymard \cite{Eymard}), and in local weighted norms $m(\cdot)$, $t(\cdot)$, $\|\cdot\|_{\mathrm{loc}}$ anchored at a link, $\|B(f,h)\|_{\mathrm{loc}}\le3(m(f)t(h)+m(h)t(f))\le\frac23\|f\|_{\mathrm{loc}}\|h\|_{\mathrm{loc}}$, by Lemma~\ref{lem:casimir}. With $q_5=\sum_{i\le5}\xi^iv_i$, the residual $\xi v_1+B(q_5,q_5)-q_5$ has exactly the degrees $6$--$10$ and local norm at most $\delta_5(\xi)$, and $J_5=2B(q_5,\cdot)$ has norm at most $\ell_5(\xi)$. For $0<\xi\le\alpha_5$ a Neumann inverse and a binary-tree majorant give $v=q_5+w$ with $\|w\|_{\mathrm{loc}}\le w_*=\frac34(1-\ell_5-\sqrt{\mathcal D_5})$, the small root of $\frac23w^2-(1-\ell_5)w+\delta_5=0$, endpoint included. The bounds are uniform in the box: each $v_n$ is a sum of universal cluster functions over plaquette multisets, so the anchored sums of the infinite lattice bound every box; the orientation dependence of the Fourier-algebra norm is handled by taking the maximum over the eight axis reflections.

\begin{proposition}[spectral return]\label{prop:return}
Under the analytic core, every eigenvalue $\lambda>0$ of $H_L-E_0$ on the physical space satisfies $\lambda\ge\kappa d_5(\xi)$.
\end{proposition}
\status{Proved here, given the analytic core; the workbench's F35--F37, audited.}
\begin{proof}
Let $\varphi$ be a physical eigenvector, $h=\varphi/\psi$ (gauge invariant) and $h'=Q_Hh\neq0$ (as $\lambda\neq0$, $h$ is not constant). By Lemma~\ref{lem:gst} and $\Gamma(v,1)=0$, $(K-\lambda/\kappa)h'=2Q_H\Gamma(v,h')$. Since $h$ is smooth, $\|Kh'\|_X=\sum_jc_j\|A_j(h')\|_1$ is finite, by Plancherel and the Cauchy--Schwarz inequality on the finite product (the workbench's F36); this is used below when dividing by it. If $\lambda/\kappa\ge3$ there is nothing to prove, since $d_5\le3$. Otherwise $c_j\ge3$ on every nonconstant physical block (Lemma~\ref{lem:casimir}), so $c_j-\lambda/\kappa\ge c_j(1-\lambda/(3\kappa))$ and $\|(K-\lambda/\kappa)h'\|_X\ge(1-\lambda/(3\kappa))\|Kh'\|_X$. The core gives $\|2\Gamma(v,h')\|_X\le4t(v)\|Kh'\|_X\le\chi\|Kh'\|_X$ with $\chi=4\sum t_i\xi^i+\frac23w_*$, and $Q_H$ does not increase $\|\cdot\|_X$. Hence $\lambda\ge3\kappa(1-\chi)$, and substituting $w_*$ shows $3(1-\chi)=d_5(\xi)$.
\end{proof}

Nothing in this argument uses the truncation order.

\begin{corollary}[the gap at every truncation order]\label{cor:ym-orderN}
For $1\le N\le5$ let $\ell_N(x)=\sum_{i\le N}(m_i+4t_i)x^i$, $\delta_N(x)=\sum_{i,j\le N,\ i+j\ge N+1}3(m_it_j+m_jt_i)x^{i+j}$, $\mathcal D_N=(1-\ell_N)^2-\frac83\delta_N$, $\alpha_N$ the first positive root of $\mathcal D_N$ and $d_N(x)=\frac32(1+\sqrt{\mathcal D_N(x)})+\sum_{i\le N}(\frac32m_i-6t_i)x^i$. If $(m_i,t_i)$, $i\le N$, bound the local norms of the first $N$ vacuum coefficients, then $\Delta_L\ge\kappa\,d_N(\xi)$ for every $L\ge2$, $a>0$ and $0<\xi\le\alpha_N$, and $d_N$ decreases on $[0,\alpha_N]$. In particular:
\begin{enumerate}[label=(\alph*)]
\item ($N=1$) $\mathcal D_1(\xi)=1-256\xi/3$ exactly and $\frac32m_1-6t_1=0$, so
\[
\Delta_L\ge\tfrac32\kappa\Bigl(1+\sqrt{1-\tfrac{64}{3g^4}}\Bigr)\qquad\text{for all }g^2\ge8/\sqrt3\approx4.6188;
\]
so $\Delta_L\ge\frac32\kappa$ there, with $2.0745\kappa$ at $g^2=5$ and $2.8304\kappa$ at $g^2=10$.
\item ($N=2$) With the workbench's pair $(m_2,t_2)=(\frac{5834}{39},\frac{137}6)$: $\alpha_2=0.0154800849822\ldots$, that is $g^2\ge4.0186791538833\ldots$, and $\Delta_L\ge\kappa\,d_2(\xi)\ge\kappa\,d_2(\alpha_2)=1.52094\,\kappa$ on that range ($1.7665\kappa$ at $g^2=4.1$, $2.3033\kappa$ at $g^2=5$). With the first pass's recomputed values, rounded up to $(m_2,t_2)=(134.0674,19.7954)$: $\alpha_2=0.0157977\ldots$, that is $g^2\ge3.97807\ldots$, with $\Delta_L\ge1.52054\,\kappa$ on that range and $\Delta_L\ge1.64707\,\kappa$ at $g^2=4$.
\item ($N=4$) $\alpha_4=0.0181049722316\ldots$, that is $g^2\ge3.7159603625\ldots$, with $d_4(1/64)=1.898118$ and $d_4(4/225)=1.658436$.
\end{enumerate}
\end{corollary}
\status{Proved here from the audited core. (a) is the workbench's own order-1 theorem, in its 14--16~September continuation YM-06 (\texttt{consolidation/20260916}, equations G22 and G26--G29, with $\Delta_L>2.07445\kappa$ at $g^2\ge5$), whose own proof was not read; its input, order~1, is verified by hand. (b) is \textsf{Generalised here} relative to the workbench's single-norm order-2 theorem in the same continuation (R10: $\Delta_L\ge\frac32\kappa\bigl(1+\sqrt{1-\frac{256}3\xi+\frac{10720}9\xi^2}\bigr)$ for $g^2\ge\sqrt{(32+\sqrt{354})/3}\approx4.1156$). With the workbench's two-norm pair and no further input, the threshold falls to $4.0187$, and the endpoint constant rises from $\frac32$ to $1.52094$. The first-pass audit (\note{46a\_}) computed $m(v_2)=134.0673186784$ and $t(v_2)=19.7953312398$ from explicit tensors. The referee pass of this paper found the closed forms $m(v_2)=6+\frac{11200}{117}+\frac{112}9+\frac{448\sqrt3}{39}$ and $t(v_2)=\frac32+\frac{80}9+\frac{256}{39}+\frac{64\sqrt3}{39}$, which agree with them to ten digits. Rounded up (with $\sqrt3<1.7320509$), they give the second threshold $3.9781$ (recomputed here by root isolation). The same continuation reaches $g^2\ge3.8260$ at order~2 (its L17) by a sharper computed bound for the cubic term of the residual ($944984/351$ in its L10, against the generic $6(m_1t_2+m_2t_1)=300672/39$), which was not audited. So the qualitative content of Theorem~\ref{thm:ym01} holds with audited inputs for $g^2\ge3.9781$, which includes the workbench's benchmark $g^2=4$, and orders 3--5 lower the threshold to $3.6836$ and raise the constants. (c) is conditional on the order-3 and order-4 inputs and does not use order~5. $\alpha_N$, the thresholds and the values of $d_N$ computed here from the exact rationals of F26 by root isolation; the monotonicity checked on a $401$-point grid for $N=1,\dots,5$. The workbench's YM-03 edition uses the radius $3/256=\alpha_1$ for its row remainder.}
\begin{proof}
Replace $q_5$ by $q_N=\sum_{i\le N}\xi^iv_i$ in the analytic core. By $v_n=\sum_{i+j=n}B(v_i,v_j)$ for $2\le n\le N$, the residual is $\xi v_1+B(q_N,q_N)-q_N=\sum_{i,j\le N,\,i+j\ge N+1}\xi^{i+j}B(v_i,v_j)$, of local norm at most $\delta_N(\xi)$, and $\|2B(q_N,\cdot)\|\le\ell_N(\xi)$. On $[0,\alpha_N]$, $\ell_N<1$: $\ell_N(0)=0$, and $\ell_N(\xi_1)=1$ with $0<\xi_1\le\alpha_N$ would give $\mathcal D_N(\xi_1)=-\frac83\delta_N(\xi_1)<0$. So the Neumann inverse and the tree majorant give $\|w\|_{\mathrm{loc}}\le w_*=\frac34(1-\ell_N-\sqrt{\mathcal D_N})$, and Proposition~\ref{prop:return} with $\chi=4\sum_{i\le N}t_i\xi^i+\frac23w_*$ gives $\lambda\ge3\kappa(1-\chi)=\kappa\,d_N(\xi)$. Differentiating $\frac23w_*^2-(1-\ell_N)w_*+\delta_N=0$ gives $w_*'\sqrt{\mathcal D_N}=\delta_N'+\ell_N'w_*>0$ on $(0,\alpha_N)$, so $\chi$ increases and $d_N$ decreases. For $N=1$, $(m_1,t_1)=(\frac{64}3,\frac{16}3)$ gives $\ell_1=\frac{128}3\xi$, $\delta_1=6m_1t_1\xi^2=\frac{2048}3\xi^2$ and $\mathcal D_1=1-\frac{256}3\xi+(\frac{16384}9-\frac{16384}9)\xi^2$; with $\xi=1/(4g^4)$, $\frac{256}3\xi=\frac{64}{3g^4}$. For $N=2$, $\frac32m_2-6t_2=\frac{3408}{39}>0$, so $d_2>\frac32$ on the whole range even before the value $d_2(\alpha_2)$ is computed.
\end{proof}

The same argument gives three further statements. They were found by the referee pass of version~1; the proofs below were checked here.

\begin{corollary}[the gap on the full space]\label{cor:fullspace}
Under the hypotheses of Corollary~\ref{cor:ym-orderN}, every eigenvalue $\lambda>0$ of $H_L-E_0$ on the whole space $L^2$, not only on the physical space, satisfies $\lambda\ge\frac14\kappa\,d_N(\xi)$. This is sharp at $\xi=0$, where the gap on $L^2$ is $\frac34\kappa$: one link in spin $\frac12$.
\end{corollary}
\begin{proof}
The ground state $e^v/\|e^v\|$ is also the ground state on $L^2$: it is positive, and the positive ground state of $H_L$ is unique. Lemma~\ref{lem:gst} holds for every $h$, gauge invariant or not. The bound $\|2\Gamma(v,h)\|_X\le4t(v)\|Kh\|_X$ of the core (the workbench's F10) uses only $\sum_ej_e\le\frac23c_j$ on the blocks of $h$, which holds on every block because $j(j+1)\ge\frac32j$ for $j\ge\frac12$. On $L^2$ the nonconstant blocks have $c_j\ge\frac34$ instead of $3$, so the proof of Proposition~\ref{prop:return} gives $\lambda\ge\frac34\kappa(1-\chi)=\frac14\kappa\,d_N(\xi)$.
\end{proof}
\status{Proved here. The record states the case $N=1$ (web continuation, (G27), p.~66).}

\begin{proposition}[the gap at complex coupling]\label{prop:complexgap}
Suppose the inputs of Corollary~\ref{cor:ym-orderN} are valid through order $N$, and let $\zeta\in\C$ with $|\zeta|\le\alpha_N$. For every $L\ge2$, on the physical space:
\begin{enumerate}[label=(\alph*)]
\item $H_L(\zeta)e^{v(\zeta)}=E_0(\zeta)e^{v(\zeta)}$ with $E_0(\zeta)=\kappa\bigl(2M_L\zeta-P_H\Gamma(v,v)\bigr)$, which is analytic in $|\zeta|<\alpha_N$;
\item every other eigenvalue $E_0(\zeta)+\lambda$ of $H_L(\zeta)$ has $\re\lambda\ge\kappa\,d_N(|\zeta|)$;
\item $E_0(\zeta)$ is an algebraically simple eigenvalue.
\end{enumerate}
Hence the ground-state Riesz projection of $H_L(\zeta)$ is analytic in $|\zeta|<\alpha_N$, for every $L$. By contrast, the global route of Proposition~\ref{prop:neumann} certifies only $|\zeta|<3/(4M_L)$, which is at most $1/320$.
\end{proposition}
\begin{proof}
(a) The fifth-reference manuscript states the core for complex $\xi$ with $|\xi|\le\alpha$ (its text after (F29), and (F31) with $x=|\xi|$), and that the same exponential solves the complex equation (after (F34)). Lemma~\ref{lem:gst} holds for complex $v$. The coefficients $v_n$ and the correction $w$ are analytic in $\zeta$, being locally uniform limits of polynomials.
(b) Let $\varphi$ be an eigenvector with eigenvalue $E_0+\lambda$, $\lambda\neq0$, and $h'=Q_H(e^{-v}\varphi)\neq0$. As in Proposition~\ref{prop:return}, $(K-\lambda/\kappa)h'=2Q_H\Gamma(v,h')$. For every block with $c_j\ge3$ and $\re\lambda\ge0$, $|c_j-\lambda/\kappa|\ge c_j-\re\lambda/\kappa\ge c_j(1-\re\lambda/(3\kappa))$. If $\re\lambda<0$, then $|c_j-\lambda/\kappa|\ge c_j$. With $\|2\Gamma(v,h')\|_X\le\chi(|\zeta|)\|Kh'\|_X$ and $\chi<1$, the case $\re\lambda<0$ is impossible. Otherwise $\re\lambda\ge3\kappa(1-\chi)=\kappa\,d_N(|\zeta|)$.
(c) The argument of (b) with $\lambda=0$ gives $h'=0$, so the eigenspace of $E_0$ is spanned by $e^v$. A generalised eigenvector $e^vh$ with $(H_L-E_0)(e^vh)=e^v$ would satisfy $\kappa(Kh-2\Gamma(v,h))=1$. Applying $Q_H$ gives $Kh'=2Q_H\Gamma(v,h')$, so $\|Kh'\|_X\le\chi\|Kh'\|_X$ and $h'=0$. Then $h$ is constant, $Kh=0$ and $\Gamma(v,h)=0$, and the equation reads $0=1$.
The family $H_L(\zeta)=\kappa K+\kappa\zeta(2M_L-S)$ is holomorphic of type (A) with compact resolvent, since $S$ is bounded. The simple eigenvalue $E_0(\zeta)$ is separated from the rest of the spectrum by the strip $0<\re(z-E_0)<\kappa d_N(|\zeta|)$, so its Riesz projection is analytic (Kato \cite{Kato}, Ch.~VII).
\end{proof}
\status{Proved here, given the core; its complex form is the workbench's own statement.}

\emph{What the theorem is and is not.} It is a finite-regulator, strong-coupling statement: fixed spacing, any box, bare coupling above a threshold. The continuum limit needs $a\to0$ with $g\to0$, the opposite regime, and the workbench records that its standing continuum path eventually leaves the domain (Proposition~\ref{prop:path}); the workbench claims no continuum statement, and Section~\ref{sec:directions} says which ingredients of the proof would have to be extended toward weak coupling. Mass gaps at strong coupling are classical for lattice gauge theories in the Euclidean setting (Osterwalder--Seiler \cite{OsterwalderSeiler}; cited, not re-read); the workbench's statement is an explicit bound, uniform in the box, on the whole physical space of the Kogut--Susskind $SU(2)$ Hamiltonian on open boxes, with a closed coupling threshold; whether a bound of this form is in the literature was not searched. If $m_5$ and $t_5$ were both underestimated by a factor $2$, the threshold would move to $g^2\approx3.737$ and $d_5(1/64)$ to $1.8946$; by Corollary~\ref{cor:ym-orderN}(b), errors of any size in the order-3 to order-5 inputs leave a gap above $1.52\kappa$ for $g^2\ge3.9781$. The zeta, Jacobian and $S^6$ material does not enter the proof of this theorem.


\subsection{The coupling domain and the continuum path}

\begin{proposition}[the standing continuum path leaves the domain]\label{prop:path}
Write $c=1/g^2$. Then $\xi=c^2/4$, so $\xi\le\alpha_5$ exactly when $c\le2\sqrt{\alpha_5}=0.27147\ldots$, and $|\xi|<1/55$ exactly when $c<2/\sqrt{55}=0.26968\ldots$. On the workbench's standing path $a_n=a_02^{-n}$, $c_n=g_0^{-2}+\beta n\log2$ with $\beta>0$:
\begin{enumerate}[label=(\alph*)]
\item the first condition fails for every $n>(2\sqrt{\alpha_5}-g_0^{-2})/(\beta\log2)$, and the second for every $n\ge(2/\sqrt{55}-g_0^{-2})/(\beta\log2)$;
\item so at most $1+\lfloor(2\sqrt{\alpha_5}-g_0^{-2})/(\beta\log2)\rfloor$ points of the path lie in the domain $\xi\le\alpha_5$. Whatever $a_0$ and $g_0$, every point in the domain has spacing $a_n\ge a_0e^{-2\sqrt{\alpha_5}/\beta}$.
\end{enumerate}
\end{proposition}
\status{Proved here; the workbench states (a) (§H7, and the fifth-reference \texttt{README.md}); (b) was found by the referee pass of version~1.}
\begin{proof}
$\xi=1/(4g^4)=c^2/4$, and $c\mapsto c^2/4$ is increasing on $c>0$. Along the path $c_n$ increases linearly without bound. For (b), a point in the domain has $n\le(2\sqrt{\alpha_5}-g_0^{-2})/(\beta\log2)\le2\sqrt{\alpha_5}/(\beta\log2)$, so $a_n=a_02^{-n}\ge a_0e^{-2\sqrt{\alpha_5}/\beta}$.
\end{proof}
So Theorem~\ref{thm:ym01} and the heat results of Section~\ref{sec:heat} hold along the standing path only down to a spacing bounded below in terms of $a_0$ and $\beta$, as the workbench says. This limits what the results are about; it does not affect them at a fixed regulator.
